% ============================================================
%  Chapter 20 — Programmable Physics: Minimal Programmable
%               Physics Kernel (MPPK) & Law-Discovery Engine
% ============================================================
\chapter{Programmable Physics}
\label{ch:programmable-physics}

\begin{center}
\textit{%
``The laws of physics are not given \emph{a~priori}---they are the
attractors of a dynamical selection process over rule-spaces.''
}
\end{center}

\medskip\noindent
This chapter formalises the \textbf{Minimal Programmable Physics Kernel
(MPPK)}: a self-evolving computational framework that generates, executes,
evaluates, and selects candidate physical theories as emergent stable
structures over graph dynamics.  Unlike preceding chapters that \emph{derive}
physics from the CIIR axioms, the MPPK \emph{discovers} physics by
stability-selection---treating laws as attractors in rule space.

% ============================================================
\section{Core System \texorpdfstring{$\mathcal{P}=(S,R,U)$}{P=(S,R,U)}}
\label{sec:mppk-core}

\begin{definition}[MPPK System]\label{def:mppk}
A \emph{Minimal Programmable Physics Kernel} is a triple
$\mathcal{P} = (S,R,U)$ where:
\begin{enumerate}[nosep]
  \item $S = (G,X)$ with graph $G=(V,E)$, node states
        $x_i\in\mathbb{R}^d$, edge weights $w_{ij}\in\mathbb{R}$.
  \item $R$ is a set of local interaction rules (bounded nonlinear
        programs).
  \item $U\colon S\times R\to S$ is the update operator.
\end{enumerate}
\end{definition}

% ------------------------------------------------------------
\subsection{Graph State Space}

The state space is a weighted directed graph
\[
  G = (V, E, w), \qquad
  V = \{1,\dots,N\},\quad
  E \subseteq V\times V,\quad
  w\colon E\to\mathbb{R}.
\]
Each vertex carries a real vector $x_i(t)\in\mathbb{R}^d$.
The full state at time $t$ is the matrix
$X(t) = [x_1(t),\dots,x_N(t)]^\top \in \mathbb{R}^{N\times d}$.

% ------------------------------------------------------------
\subsection{Initial Kernel}

The \emph{initial kernel} (to be validated before mutation) implements:

\begin{equation}\label{eq:mppk-kernel}
  x_i(t{+}1) = x_i(t)
  + \sum_{j\in\mathcal{N}(i)} w_{ij}\,
    \tanh\!\bigl(x_j(t) - x_i(t)\bigr),
\end{equation}

where $\mathcal{N}(i) = \{j : (i,j)\in E\}$.

\begin{proposition}[Boundedness of Tanh Kernel]\label{prop:tanh-bounded}
For any graph $G$ with weights $|w_{ij}|\le W$ and maximum degree
$\Delta$, the update~\eqref{eq:mppk-kernel} satisfies
\[
  \|x_i(t{+}1) - x_i(t)\|_\infty \le W\Delta.
\]
\end{proposition}

% ============================================================
\section{Observables}
\label{sec:mppk-observables}

\begin{definition}[Energy Functional]\label{def:mppk-energy}
\begin{equation}
  E(t) = \sum_{(i,j)\in E} w_{ij}\,\|x_i(t) - x_j(t)\|^2.
\end{equation}
\end{definition}

\begin{definition}[Momentum Observable]\label{def:mppk-momentum}
\begin{equation}
  P(t) = \sum_{i\in V} \bigl(x_i(t) - x_i(t{-}1)\bigr).
\end{equation}
\end{definition}

\noindent
$E(t)$ measures total ``tension'' in the graph; $P(t)$ measures
net displacement.  Neither is \emph{assumed} conserved---conservation
is an emergent property to be \emph{mined}.

% ============================================================
\section{Six-Step Execution Loop}
\label{sec:mppk-loop}

Each generation executes the following pipeline.

% ------------------------------------------------------------
\subsection{Step 1: Simulation}

Run the current kernel for $T$ timesteps, recording the full
trajectory $\{X(0), X(1), \dots, X(T)\}$ together with
$\{E(t)\}$ and $\{P(t)\}$.

% ------------------------------------------------------------
\subsection{Step 2: Structure Extraction}

Detect emergent dynamical structures:
\begin{itemize}[nosep]
  \item \textbf{Fixed points:}
        $\|X(T) - X(T{-}1)\| < \epsilon_{\mathrm{fp}}$.
  \item \textbf{Limit cycles:}
        $\exists\,\tau:\|X(T) - X(T{-}\tau)\| < \epsilon_{\mathrm{cyc}}$.
  \item \textbf{Chaotic regimes:}
        positive Lyapunov exponent proxy $\lambda>0$.
  \item \textbf{Cluster formation:}
        $k$-cluster decomposition of $\{x_i(T)\}$.
  \item \textbf{Symmetry breaking:}
        $k>1$ in above cluster analysis.
\end{itemize}

% ------------------------------------------------------------
\subsection{Step 3: Invariant Mining}

\begin{enumerate}[nosep]
  \item \textbf{Conserved quantities:} compute relative variation
        $\sigma_Q / |\mu_Q|$ for each candidate $Q\in\{E, \|P\|, \sum\|x_i\|\}$.
        If $<\tau_{\mathrm{cons}}$ the quantity is approximately conserved.
  \item \textbf{Symmetry transformations:} detected via near-invariance
        under node permutations.
  \item \textbf{Spectral stability:} eigenvalues of the graph Laplacian
        $L = D - W_{\mathrm{sym}}$.  A positive spectral gap
        $\lambda_2(L)>0$ indicates connectivity and mixing.
\end{enumerate}

% ------------------------------------------------------------
\subsection{Step 4: Fitness Evaluation}

\begin{equation}\label{eq:mppk-fitness}
  F(\mathcal{P}) =
    \alpha\cdot\mathrm{stability}
  + \beta\cdot\mathrm{structure}
  + \gamma\cdot\mathrm{boundedness}
  + \delta\cdot\mathrm{compressibility},
\end{equation}
with $\alpha+\beta+\gamma+\delta=1$.

\begin{itemize}[nosep]
  \item \textbf{Stability:} energy conservation $+$ fixed-point/cycle
        detection $+$ spectral gap.
  \item \textbf{Structure formation:} cluster count, symmetry breaking,
        absence of chaos.
  \item \textbf{Boundedness:} $\max_t\sum_i\|x_i(t)\| < M$.
  \item \textbf{Compressibility:} $1/(1+\#\text{params}/10)$.
\end{itemize}

% ------------------------------------------------------------
\subsection{Step 5: Rule Mutation}

Five mutation strategies:
\begin{enumerate}[nosep]
  \item \textbf{Modify interaction function:} replace $\tanh$ with
        $\sigma$, quadratic, ReLU-clip, or cubic.
  \item \textbf{Alter graph topology:} rewire edges with probability $p$.
  \item \textbf{Introduce stochasticity:}
        $x_i \leftarrow x_i + \sigma\eta$, $\eta\sim\mathcal{N}(0,I)$.
  \item \textbf{Nonlinear coupling:} apply $d\times d$ coupling matrix $C$
        before the interaction function.
  \item \textbf{Perturb edge weights:}
        $w_{ij} \leftarrow w_{ij} + \epsilon$, $\epsilon\sim\mathcal{N}(0,\sigma_w^2)$.
\end{enumerate}

\noindent\textbf{Hard constraint:} reject any mutant causing
$\max_t\sum_i\|x_i(t)\| > 10^6$ (immediate divergence).

% ------------------------------------------------------------
\subsection{Step 6: Selection}

Retain the top-$k$ rule systems by fitness:
\[
  R_{t+1} = \operatorname*{arg\,max}_k\, F(R_i).
\]

% ============================================================
\section{Meta-Learning Layer}
\label{sec:mppk-meta}

After $G$ generations, the meta-learning layer infers:

\begin{enumerate}[nosep]
  \item \textbf{Effective emergent force laws:} the interaction type
        and parameters of the best-fitness rule.
  \item \textbf{Stable dynamical invariants:} quantities conserved
        in $\ge G/2$ generations.
  \item \textbf{Low-dimensional embeddings:} PCA of late-trajectory
        states yields effective dimension $d_{\mathrm{eff}}$.
  \item \textbf{Candidate physical laws:} compact verbal descriptions
        of emergent regularities.
\end{enumerate}

\begin{remark}[No hard-coded physics]
The MPPK never assumes Newton, quantum mechanics, or general
relativity.  All candidate laws are \emph{emergent attractors}
in rule space, discovered via the selection loop.
\end{remark}

% ============================================================
\section{Optional Extensions}
\label{sec:mppk-extensions}

If a stable regime emerges, the MPPK may be extended with:

\begin{description}[nosep]
  \item[Phase variables:]
    $x_i \to (x_i, \theta_i)$ with angular dynamics.
  \item[Complex-valued states:]
    $x_i \in \mathbb{C}^d$ for quantum-like interference.
  \item[Adaptive rewiring:]
    $E(t{+}1) = f\bigl(E(t), X(t)\bigr)$---edges rewired
    based on state similarity.
  \item[Multi-population co-evolution:]
    multiple competing rule-sets evolving in parallel.
\end{description}

% ============================================================
\section{Computational Implementation}
\label{sec:mppk-code}

The MPPK is implemented in
\texttt{quasim/ciir/swarm/mppk.py}
with 64 passing tests in
\texttt{tests/test\_mppk.py}.
Key classes:

\begin{center}
\begin{tabular}{lll}
\toprule
\textbf{Class} & \textbf{Role} & \textbf{Section} \\
\midrule
\texttt{MPPKGraph}      & State space $S$       & \S\ref{sec:mppk-core} \\
\texttt{MPPKRule}        & Interaction rule $R$  & \S\ref{sec:mppk-core} \\
\texttt{UpdateOperator}  & Update operator $U$   & \S\ref{sec:mppk-core} \\
\texttt{simulate()}      & Step 1                & \S\ref{sec:mppk-loop} \\
\texttt{extract\_structure()} & Step 2           & \S\ref{sec:mppk-loop} \\
\texttt{mine\_invariants()}   & Step 3           & \S\ref{sec:mppk-loop} \\
\texttt{evaluate\_fitness()}  & Step 4           & \S\ref{sec:mppk-loop} \\
\texttt{RuleMutator}     & Step 5                & \S\ref{sec:mppk-loop} \\
\texttt{MPPKEngine}      & Steps 1--6 + Selection & \S\ref{sec:mppk-loop} \\
\texttt{meta\_learn()}   & Meta-learning layer   & \S\ref{sec:mppk-meta} \\
\bottomrule
\end{tabular}
\end{center}

% ============================================================
\section{Hard Constraints}
\label{sec:mppk-constraints}

\begin{enumerate}[nosep]
  \item All interaction functions $f\colon\mathbb{R}^d\to\mathbb{R}^d$
        are bounded: $\|f(\Delta x)\|\le B$ for some finite $B$.
  \item No physical constants ($G$, $\hbar$, $c$, etc.) are hard-coded.
  \item Laws are emergent attractors in rule space, not assumed.
  \item Divergent candidates are rejected at Step~6.
  \item All primitives are explicitly defined.
\end{enumerate}

% ============================================================
\section{Terminal Objective}
\label{sec:mppk-terminal}

The MPPK evolves toward:

\begin{quote}
\emph{A self-organizing space of candidate physical laws generated by
stability-selection over graph dynamics.}
\end{quote}

\noindent
This completes the programmable physics layer of the CIIR framework,
providing a concrete computational substrate for the discovery of
law-like behavior from first principles of dynamical consistency.
